\section{代码与数据说明}

本研究的数据预处理、局地化集合调整卡尔曼滤波、交叉熵航线规划器、联合嵌入预测架构时空模型、出发前预演与在线留出估值控制器、论文全部图表的生成脚本、原始实验结果数据，以及作者与人工智能助手的研究讨论记录（已隐去涉及作者其他未发表工作的讨论内容）均开源于：\TBD{GitHub 仓库地址}。在代码仓库中，运行 \texttt{experiments/rev2.sh} 即可按固定随机种子完整重现正文与附录的全部数值实验，运行 \texttt{experiments/rev2\_summary.py} 与 \texttt{experiments/rev2\_tables.py} 可由原始输出日志直接汇总生成文中的全部统计置信区间与表格文件。

实验使用的两套卫星遥感数据集均下载自美国国家海洋和大气管理局海岸观测网的公开数据服务器：
\begin{enumerate}
  \item \textbf{东海多星融合叶绿素 a 逐日数据集（2018—2026年）：}融合了极轨卫星水色传感器观测并经经验正交函数缺失数据插补算法填补云区盲区，空间分辨率为$2\text{ km}$，时间分辨率为逐日，遵循公共领域开放许可，数据集标识为 \texttt{noaacwNPPN20S3ASCIDINEOF2kmDaily}。
  \item \textbf{多尺度超分辨率海表温度逐日数据集（2018—2026年）：}由美国国家航空航天局喷气推进实验室制作（版本4.1），原始空间分辨率为$0.01^\circ$，时间分辨率为逐日，数据集标识为 \texttt{jplMURSST41}。
\end{enumerate}

\section{人工智能工具使用说明}

在本课程论文的研究与写作过程中，人工智能工具的使用情况与各环节具体分工如下：
\begin{enumerate}
  \item 作者本人提出了本文的核心研究问题（自主水下航行器在现场巡测中何时应当继续计算、何时应当继续观测），确立了研究思路、实验设计与写作规范，并对论文的全部内容与学术诚信承担完全责任。
  \item Claude（Anthropic）协助完成了相关文献检索与核查，与作者讨论细化了技术方案，编写了仿真与实验代码，运行了全部数值实验并绘制了论文图表。
  \item Gemini（Google）根据实验事实清单与结构大纲协助撰写了论文初稿与修改稿的正文文字。
  \item 作者与 Claude 对照原始实验日志与出版文献逐一核查了文中的全部数值结果、置信区间与参考文献引用。
\end{enumerate}
